% SPDX-License-Identifier: LPPL-1.3c
% Copyright (C) 2026 Christophe Jorssen
% Author and Current Maintainer: Christophe Jorssen <christophe.jorssen@gmail.com>
% LPPL maintenance status: maintained
% This file is part of luafemm. See LICENSE and LICENSES.md for its terms.
\section{A tutorial: drawing electric fields from charges}
\label{sec:tutorial-electric}

Kamil would like to compare the electric field of one charge with that of a
dipole. He needs both a diagram and a graph, and he wants the numbers in the
graph to come from the sources in his diagram. He starts with a model whose
solution is known: prescribed point charges in vacuum.

Here the word ``point'' has its three-dimensional meaning. A charge is measured
in coulombs, and its field decreases as the inverse square of distance.
The drawing is a section through the plane containing all the charges.
It is not a two-dimensional finite-element calculation, and the charges are
neither infinite charged wires nor rings.

\luafemmexamplebegin{electric-charges}
\subsection{One node, one physical charge}
The usual |\usetikzlibrary{femm}| loads the electric interface as well as the
magnetic one. Kamil also loads |backgrounds| so that he can draw the field behind
symbols that have already declared their physical sources.
He puts |femm/electrostatic| on the picture and a |femm/point charge| on a node.
The source is at the node's \emph{centre}, regardless of its label.

\begin{codeexample}[setup code]
\usetikzlibrary{backgrounds}
\end{codeexample}

\begin{codeexample}[width=0pt]
\begin{tikzpicture}[femm/electrostatic={xmin=-2,xmax=2,ymin=-2,ymax=2}]
  \node[femm/point charge={charge=1e-9},fill=red!15] at (0,0) {$+$};
  \begin{scope}[on background layer]
    \pic[femm/electric={levels={600,900,1800}}] {femm equipotentials};
    \pic {femm electric field};
  \end{scope}
\end{tikzpicture}
\end{codeexample}

The charge is $1\,\mathrm{nC}$. By default a coordinate unit represents one
centimetre physically, and TikZ displays it as one centimetre on the page.
The blue curves follow $\mathbf E$; the dashed curves have constant electric
potential. Their specified levels are in volts. In vacuum this charge has
$V\simeq899\,\mathrm V$ at a distance of one centimetre.

The small circle is a \emph{symbol}, not the surface of a charged sphere.
|minimum size| changes the symbol; |exclusion radius| changes the disk omitted
from plots near the singular source. Neither changes Coulomb's law. Its default
radius is $0.22$ coordinate units, convenient for a four-millimetre symbol at the
default scale. At a different drawing scale Kamil can adjust these two controls
independently.

There is no physical mesh to refine in this example. Kamil first checks his
charges and units. He can use |contour resolution=30| while arranging the figure,
then increase it to 80 or 160 for the final equipotentials. This refines the
\emph{drawing}, not the exact analytical values used by the graphs.

\subsection{Two charges and a contour to measure along}
Kamil now makes a dipole. Reusable TikZ styles keep the charge definitions and
symbol colours together. He draws the measurement contour one centimetre above
the charges, where the field is regular.

\begin{codeexample}[setup code]
\tikzset{
  positive charge/.style={femm/point charge={charge=1e-9},fill=red!15},
  negative charge/.style={femm/point charge={charge=-1e-9},fill=blue!15}}
\end{codeexample}

\begin{codeexample}[width=0pt]
\begin{tikzpicture}[femm/electrostatic={name=lesson dipole}]
  \node[positive charge] (positive) at (-1,0) {$+$};
  \node[negative charge] (negative) at (1,0) {$-$};
  \begin{scope}[on background layer]
    \pic[femm/electric={levels={-900,-450,0,450,900}}]
      {femm equipotentials};
    \pic[femm/electric={lines per charge=12}] {femm electric field};
  \end{scope}
  \draw[thick,densely dotted,
    femm/electric profile={name=above dipole,samples=161}]
    (-2,1)--(2,1) node[right] {$\mathcal C$};
\end{tikzpicture}
\end{codeexample}

The order matters. First declare \emph{all} the source nodes, then request
fields or profiles. The first evaluation freezes the source list. Adding another
physical charge afterwards is an error: otherwise an earlier field picture and
a later graph could describe different problems. Ordinary labels, arrows and
other annotations can still be added at any time.

The profile has a document-wide name. It owns its source data and remains
available after the picture ends. With PGFPlots loaded, Kamil can now plot the
potential without reproducing a single charge definition:

\begin{codeexample}[width=0pt]
\begin{tikzpicture}
  \begin{axis}[width=9cm,height=5cm,no markers,grid=major,
    xlabel={$x$ (cm)},ylabel={$V$ (V)}]
    \femmelectricplot[profile=above dipole,component=V,abscissa=x,
      plot options={blue!70!black,thick}]
  \end{axis}
\end{tikzpicture}
\end{codeexample}

The potential vanishes on the perpendicular bisector. That does not mean the
field vanishes there: the two horizontal field components reinforce each other.
Kamil uses the same sampled contour to display both components and their norm.

\begin{codeexample}[width=0pt]
\begin{tikzpicture}
  \begin{axis}[width=9cm,height=5cm,no markers,grid=major,
    xlabel={$x$ (cm)},ylabel={Electric field (V/m)},
    legend style={at={(.5,1.03)},anchor=south,legend columns=3}]
    \femmelectricplot[profile=above dipole,component=Ex,abscissa=x,
      plot options={blue!70!black,thick}]
    \addlegendentry{$E_x$}
    \femmelectricplot[profile=above dipole,component=Ey,abscissa=x,
      plot options={red!70!black,dashed,thick}]
    \addlegendentry{$E_y$}
    \femmelectricplot[profile=above dipole,component=E,abscissa=x,
      plot options={black,densely dotted,thick}]
    \addlegendentry{$|\mathbf E|$}
  \end{axis}
\end{tikzpicture}
\end{codeexample}

\subsection{Keeping the model, changing the view}
The picture's |name=lesson dipole| saved an independent source snapshot at the
end of that picture. Kamil reuses it for a vector view. Each arrow below has the
same length, so it indicates direction only. A quantitative arrow map instead
uses |vector mode=scaled| and a chosen |vector scale|.

\begin{codeexample}[width=0pt]
\begin{tikzpicture}[femm/electrostatic={snapshot=lesson dipole}]
  \node[circle,draw,fill=red!15] at (-1,0) {$+$};
  \node[circle,draw,fill=blue!15] at (1,0) {$-$};
  \pic[femm/electric={vector spacing=.4,vector length=.22}]
    {femm electric vectors};
\end{tikzpicture}
\end{codeexample}

These are ordinary circles used as annotations. Repeating physical charge nodes
in a restored snapshot would try to add sources to an already frozen model and
is rejected. The snapshot stores source coordinates and physical units; its new
picture supplies the display transform. Thus |scale|, |rotate|, and TikZ's |x|
and |y| vectors can change the appearance without changing the calculation.

\subsection{Choosing which lines to show}
Automatic lines start just outside positive charge masks, or negative masks if
there are no positive charges. The line count is a drawing choice; it is not a
measurement of three-dimensional flux. In particular, a net negative collection
can have incoming lines from infinity that positive-source seeding will miss.
Kamil can set |seed sign=negative| or give his own starting points.

A TikZ contour is often more convenient than a list of coordinates. This one
places nine seeds uniformly in arclength along a short segment. Each line is
integrated in both directions and its arrows always point along $\mathbf E$.

\begin{codeexample}[width=0pt]
\begin{tikzpicture}[femm/electrostatic={snapshot=lesson dipole}]
  \node[circle,draw,fill=red!15] at (-1,0) {$+$};
  \node[circle,draw,fill=blue!15] at (1,0) {$-$};
  \draw[gray,densely dotted,femm/electric seeds={samples=9}]
    (0,.2)--(0,1.8);
  \pic[femm/electric={step=.05,tolerance=.00001}] {femm electric field};
\end{tikzpicture}
\end{codeexample}

\subsection{More charges, the same declarations}
A quadrupole requires no new solver settings. Local TikZ transformations position
the source centres just as they position the visible nodes. The empty-looking
central region contains a null point, where a field direction is undefined;
line integration stops if it reaches such a point.

\begin{codeexample}[width=0pt]
\begin{tikzpicture}[femm/electrostatic={xmin=-2.5,xmax=2.5},scale=.9]
  \foreach \x/\y in {-1/-1,1/1}
    \node[positive charge] at (\x,\y) {$+$};
  \foreach \x/\y in {-1/1,1/-1}
    \node[negative charge] at (\x,\y) {$-$};
  \begin{scope}[on background layer]
    \pic[femm/electric={levels={-700,-300,300,700}}] {femm equipotentials};
    \pic[femm/electric={lines per charge=12}] {femm electric field};
  \end{scope}
\end{tikzpicture}
\end{codeexample}

Finally Kamil tries an asymmetric collection. The charge values remain explicit;
the text inside a node does not set its charge. Superposition works for any
finite number of prescribed coplanar charges in the same homogeneous medium.
Putting a conducting outline or a dielectric region in the drawing would not
redistribute those charges or change the field.

\begin{codeexample}[width=0pt]
\begin{tikzpicture}[femm/electrostatic={relative permittivity=2,
    xmin=-3,xmax=3,ymin=-2,ymax=2}]
  \node[femm/point charge={charge=2e-9},fill=red!15] at (-1.4,-.4) {$2q$};
  \node[femm/point charge={charge=-1e-9},fill=blue!15] at (.9,.7) {$-q$};
  \node[femm/point charge={charge=-.5e-9},fill=blue!10] at (.6,-1) {$-q/2$};
  \begin{scope}[on background layer]
    \pic[femm/electric={levels={-300,0,300,600,900}}] {femm equipotentials};
    \pic[femm/electric={lines per charge=20}] {femm electric field};
  \end{scope}
\end{tikzpicture}
\end{codeexample}

Here $q=1\,\mathrm{nC}$ and $\epsilon_r=2$. Multiplying a uniform permittivity
changes the magnitudes of $V$ and $\mathbf E$, but not the direction of the lines.
Only their interaction with a nonuniform medium would require a different
physical model.
\luafemmexampleend{electric-charges}
